\clearpage
\section{Individual Contribution Statements}
\label{app:ics}
The statements below describe work at the time each member prepared their
individual record. Historical page and test counts refer to those dated
builds; the current reproduction checklist is in \path{report_v3/checklist/README.md}.
% Source: Project1/ICS/TianruiJiICS.tex
\begin{formbox}
\small
\textbf{Name:} Tianrui Ji \hfill \textbf{Student ID:} 2361367\\
\textbf{Role this project:} Algorithms, numerical experiments and visualisation.

\textbf{1. My specific contributions to this project:}

\begin{itemize}[leftmargin=1.4em,itemsep=1pt,topsep=2pt]
    \item In \path{Project1/code/model.py}, I implemented column-wise matrix
    vectorisation, the right-hand side, analytic Jacobian, finite-time SVD
    solution, energy, singular values, and orthogonality diagnostics. The
    finite-time matrix, rather than only the limiting polar factor, is the
    reference for numerical errors.
    \item In \path{Project1/code/solvers.py}, I implemented Explicit Euler,
    classical RK4, and Implicit Euler with damped Newton iteration. I also
    implemented fixed-grid and step-doubling adaptive drivers, failed-step
    handling, and work counters, then checked that the implicit residual and
    Jacobian matched the equations in Section~2.
    \item In \path{Project1/code/experiments.py} and
    \path{Project1/code/run_all.py}, I developed the reproducible experiments
    and eight baseline figures used in Sections~3--4, including convergence,
    stability, trajectories, cost--accuracy, adaptivity, and rank-deficient
    behaviour. I checked the associated CSV/JSON evidence against the exact
    solution and numerical claims.
    \item For Section~5, I implemented the separate two-stage Gauss IRK4
    method in \path{Project1/code_advanced/gauss_irk4.py}, including a coupled
    $18\times18$ damped Newton solve. I wrote the four-method comparison and
    four additional figures in \path{code_advanced/run_comparison.py}; its
    refined-grid slope is about $3.9674$, and the matched-error comparison
    records 720 versus 300 RHS evaluations for Gauss IRK4 and classical RK4.
    These counts are numerical work measures, not wall-clock timings.
\end{itemize}

\textbf{2. The hardest conceptual obstacle I encountered:}

The hardest implementation issue was solving both Gauss stages as one coupled
18-dimensional nonlinear system rather than as two independent stages. I
constructed its block residual Jacobian, used damped Newton steps that reduce
the full residual, and checked the resulting trajectory against the exact
finite-time matrix. This also forced me to distinguish a small stage residual
from a small global time-discretisation error.

\textbf{3. One piece of AI-generated output my team used, and what I changed:}

During the AI-assisted review of Newton tolerance described in Appendix~A,
a proposed explanation treated a fixed $10^{-12}$ residual as sufficient to
guarantee time-stepping accuracy or asymptotic order. I corrected the
numerical interpretation: the threshold controls the
implicit algebraic solve, whereas the final error also contains time
discretisation error. I checked the implemented residual and compared terminal
states with the exact matrix; the fixed threshold is supported over the
tested finite grids, not as an unrestricted $h\to0$ guarantee.

\textbf{4. One thing I still do not fully understand:}

I cannot yet derive a sharp global nonlinear step-size guarantee for the full
matrix trajectory from the frozen Jacobian spectrum. In particular, the
negative-real-axis amplification plot and Gauss IRK4's A-stability do not by
themselves certify nonlinear behaviour at arbitrary step sizes.

\vspace{0.5em}
\textbf{Signature:} \underline{\hspace{4.5cm}} \hfill
\textbf{Date:} \underline{\hspace{2.5cm}}
\end{formbox}
\clearpage
% Source: Project1/ICS/Jingyan.WeiCSi.tex
\begin{formbox}
\footnotesize
\textbf{Jingyan Wei} (Student ID: 2361411)\\
\textbf{Role this project:} Mathematical Theory (Schemes and Stability Lead)

\textbf{1. My specific contributions to this project:}
\begin{itemize}[leftmargin=1.3em,itemsep=1pt,topsep=1pt]
    \item In Section~2.1, I derived the updates for Explicit Euler, classical
    RK4 with its Butcher tableau, and Implicit Euler. I derived the matrix
    Fr\'echet derivative
    $\mathcal{D}f(X)[H]=H(I_n-X^TX)-X(H^TX+X^TH)$, the residual Jacobian
    $J_F=I_{n^2}-hJ_f(X)$, and the damped Newton iteration with backtracking.
    \item In Section~2.3, I derived the Dahlquist amplification functions,
    the RK4 negative-real intercept $a^*\approx2.78529$ and imaginary-axis
    limit $[-2\sqrt2,2\sqrt2]$, and verified the A- and L-stability of
    Implicit Euler.
    \item For Topic~5, I transformed the linearised operator into SVD
    coordinates, separated contracting symmetric modes from neutral
    skew-symmetric modes, computed the initial spectrum
    $\{-26.0,\ldots,+0.88\}$, obtained
    $h_{\mathrm{crit}}^{\mathrm{EE}}\approx0.0769$ and
    $h_{\mathrm{crit}}^{\mathrm{RK4}}\approx0.1071$, and wrote the disclaimer
    that these frozen-spectrum thresholds are local diagnostics.
    \item In Section~5, I extended the stability analysis to two-stage Gauss
    IRK4: I derived its rational stability function, proved A-stability from
    $|D(z)|^2-|N(z)|^2\geq0$, documented the failure of L-stability through
    $R_G(z)\to1$ as $z\to-\infty$, and derived the coupled $18\times18$ block
    Newton Jacobian.
\end{itemize}

\textbf{2. The hardest conceptual obstacle I encountered:}
The main challenge was deriving the problem-specific Fr\'echet derivative and
Jacobian spectrum for $\dot X=X(I_n-X^TX)$. Matrix multiplication is
non-commutative, so directional differentiation of the cubic term produces
three ordered terms. Diagonalising each coupled $2\times2$ coordinate block
then showed that symmetric perturbations contract while skew-symmetric
perturbations form neutral tangent rotations.

\textbf{3. One piece of AI-generated output my team used, and what I changed:}
An AI response claimed that $|1+h\lambda_j|\leq1$ for the initial frozen
spectrum proves global stability on $[0,2]$. I corrected this because
$J_f(X(t))$ changes along the nonlinear trajectory, so the spectrum is only a
local diagnostic. I added the analytical disclaimer and the empirical sweep
protocol (Appendix~A, Audit~2). I also corrected the claim that Gauss IRK4 is
L-stable (Audit~4).

\textbf{4. One thing I still do not fully understand:}
I do not yet have a sharp nonlinear perturbation bound along the entire
trajectory. Although $J_f(X)$ is symmetric at each state, its eigenvalues and
eigendirections evolve, so pointwise frozen spectra do not provide a global
time-dependent certificate over $[0,2]$.

\textbf{Signature:} \underline{\hspace{4.5cm}} \hfill
\textbf{Date:} \underline{\hspace{2.5cm}}
\end{formbox}
\clearpage
% Source: Project1/ICS/MeichenDingICS.tex
\begin{formbox}
\footnotesize
\textbf{Meichen Ding} (Student ID: 2362344)\\
\textbf{Role this project:} Mathematical Theory Lead (Consistency,
Adaptivity and Convergence)

\textbf{1. My specific contributions to this project:}
\begin{itemize}[leftmargin=1.3em,itemsep=1pt,topsep=1pt]
    \item In Section~2.2, I wrote the consistency and local truncation error
    foundations under Lipschitz continuity and the small-step condition
    $hL<1$. I verified Euler defects of order $\mathcal{O}(h^2)$ and the RK4
    defect of order $\mathcal{O}(h^5)$ using the fourth-order conditions.
    \item In Sections~2.2.3--2.2.4, I derived the step-doubling Richardson
    estimator with denominator $2^p-1$, and the scale-invariant controller
    with exponent $-1/(p+1)$, safety factor $0.9$, and limiter $[0.2,2.0]$.
    \item In Section~2.4, I established the discrete error recurrence under a
    uniform increment Lipschitz bound and bounded residual-Jacobian inverse. I
    derived the discrete Gr\"onwall bound
    $\|e_n\|\leq[(e^{\Lambda(T-t_0)}-1)/\Lambda]\tau_{\max}(h)$ and explained
    how local defects accumulate into global order $p$. I also derived the
    log--log slope relation used in the experiments.
    \item In Sections~5.3 and~5.5, I extended the convergence discussion to
    Gauss IRK4: I identified its order $p=2s=4$, stated the stage contraction
    condition $h\|A\|_\infty L<1$, and analysed how inexact stage residuals
    must remain below the time-discretisation scale on the tested grids.
\end{itemize}

\textbf{2. The hardest conceptual obstacle I encountered:}
The main challenge was explaining why accepting the two-half-step value $y_f$
does not increase the base method's global order from $p$ to $p+1$. Expanding
the coarse and fine steps showed that step doubling reduces the leading error
coefficient by $2^{-p}$ but does not satisfy additional order conditions. This
also explains the controller exponent $-1/(p+1)$.

\textbf{3. One piece of AI-generated output my team used, and what I changed:}
An AI response proposed $h_{\mathrm{new}}=hE^{-1/p}$. I identified that it
mixed global order $\mathcal{O}(h^p)$ with a local defect
$\mathcal{O}(h^{p+1})$. I corrected the exponent to $-1/(p+1)$, derived the
$2^p-1$ denominator, and added the safety factor and limiter (Appendix~A,
Audit~1). I also checked the finite-grid interpretation of the Gauss stage
residual in Section~5.

\textbf{4. One thing I still do not fully understand:}
I do not yet know how to construct a rigorous non-autonomous matrix logarithmic
norm for this nonlinear flow over $[0,2]$. Such a bound could relate global
contraction to the observed transient behaviour of the discrete Lyapunov
residual.

\textbf{Signature:} \underline{\hspace{4.5cm}} \hfill
\textbf{Date:} \underline{\hspace{2.5cm}}
\end{formbox}
\clearpage
% Source: Project1/ICS/JiashengXuICS.tex
\begin{formbox}
\textbf{Name:} Jiasheng Xu \hfill \textbf{Student ID:} 2251640\\
\textbf{Role:} Testing and validation.

\textbf{1. My specific contributions to this project:}
\begin{enumerate}[leftmargin=2em,itemsep=0.25em,topsep=0.15em]
    \item I added and migrated three edge-case validation tests in
    \path{Project1/test/test_validation.py}: tighter adaptive tolerances must
    reduce global error, an oversized explicit step exposes singular-value
    crossing, and Newton failure triggers adaptive retry.  These changes are
    traceable to commit \texttt{c4ca4ee} and merge commit \texttt{b28fe39}.

    \item I documented the earlier clean-environment validation in
    \path{Project1/report_v3/checklist/VALIDATION_RUN_2026-10-03.md}
    (commit \texttt{06461a9}).  For the final English-only submission, I
    independently reran the English entry point and the complete core test
    suite, obtaining 13/13 passing tests.  After review, I aligned the report's
    Radau comparison with its bundled data at \(1.16\times10^{-13}\) and
    regenerated and independently checked all 41 manifest hashes in the
    Windows submission checkout using the documented procedure (commit
    \texttt{838d938}).

    \item After the validation merge, I rebuilt the report PDF from
    \path{Project1/report_v3/report.tex}, checked all 43 pages, and confirmed
    that Section~4 reports thirteen core tests and that the contents, references,
    and appendices are complete.  I also updated the change log, validation
    checklist, and affected manifest hashes.  This report-integration work was
    merged through PR~\#3 (merge commit \texttt{e6913e6}).
\end{enumerate}

\textbf{2. The hardest conceptual obstacle I encountered:} The main challenge
was separating a local step-doubling error estimate from accumulated global
error.  I addressed it by comparing against the exact solution and turning the
expected global-error reduction into a deterministic regression test rather
than treating one accepted local step as evidence of global accuracy.

\textbf{3. One piece of AI-generated output my team used, and what I changed:}
As recorded in Appendix~A, AI helped check deliverable completeness and
organise a draft of the validation evidence.  I revised that draft after manual
reproduction: I corrected the Radau rounding, distinguished close agreement
from exact matching, ran the complete suite in a clean environment, and
independently verified the documented SHA-256 results before submission.

\textbf{4. One thing I still do not fully understand:} I do not yet have a
rigorous global-error bound for this nonlinear matrix flow that starts from the
local step-doubling estimator while also accounting for rejected Newton steps.

\vfill
\textbf{Signature:} \underline{\hspace{4.5cm}} \hfill
\textbf{Date:} \underline{\hspace{2.5cm}}
\end{formbox}
\clearpage
% Source: earlier report_v3 Appendix B statement supplied by Han Liu.
\begin{formbox}
\textbf{Han Liu} (Student ID: 2364393)
\par
\begingroup\small
\setlength{\parskip}{0.2em}

\textbf{1. My specific contributions to this project:}

I was mainly responsible for the final-stage writing and integration of the
report. I completed and refined Section~1 (Background) and Section~5
(Conclusions), including the project motivation, overall context, and the
synthesis of the main numerical findings. I also checked the completeness and
consistency of the final report files, updated the appendices and version
documentation, and helped integrate the existing Sections~2--4, figures, data,
and references into the complete report. In the final checking stage, I also
reviewed several wording and notation issues in Section~2 and clarified the
discussion of the Implicit Euler Newton tolerance and local error terminology.

\vspace{0.2em}
\textbf{2. The hardest conceptual obstacle I encountered:}

The most difficult part for me was understanding how the theoretical error
statements in Section~2 should be connected to the actual numerical
implementation. In particular, I had to distinguish the per-unit-step local
truncation error from the unnormalised local one-step defect, and understand
why the convergence proof for Implicit Euler assumes a sufficiently accurate
nonlinear solve. I addressed this by comparing the mathematical definitions,
the Newton stopping criterion used in the code, and the numerical convergence
results before revising the wording in the report.

\vspace{0.2em}
\textbf{3. One piece of AI-generated output my team used, and what I changed:}

As recorded in Appendix~A, AI was used to help review the report structure and
identify wording and completeness issues in the final document. One example was
the suggested clarification of the Implicit Euler convergence discussion in
Section~2. I checked the AI-assisted explanation against the existing
derivation and numerical results, and revised the wording so that the report
explicitly states that a fixed Newton tolerance is adequate for the tested
finite range of step sizes but is not itself an asymptotic assumption as
$h\to0$. The final wording and its inclusion in the report were decided after
this checking process.

\vspace{0.2em}
\textbf{4. One thing I still do not fully understand:}

I still do not fully understand the more general theory describing how the
accuracy of an inexact nonlinear solve affects the convergence order of an
implicit time-stepping method beyond this specific Implicit Euler example. I
understand the practical requirement that the nonlinear-solver error should be
smaller than the time-discretisation error, but I would like to understand a
more general rigorous analysis of this interaction.

\vspace{0.3em}
\textbf{Signature:} \underline{\hspace{4.5cm}}
\qquad
\textbf{Date:} \underline{\hspace{2.5cm}}

\par\endgroup
\end{formbox}
\clearpage
